\section{Chellas's classical systems and neighborhood constructions}
\label{sec:p08-n-b-006-03}

Chellas's Part~III develops minimal models, also called neighborhood
models, alongside classical, monotonic and regular modal systems.
Its proof methods distinguish the smallest and largest canonical
models, supplementation, finite intersection closure and augmentation
\cite[chaps.~7--9, pp.~207--267]{Chellas1980}.
This presentation supplies those constructions for all eight combinations
of the source's schemas M, C and N. The source's names \(M\) and \(R\)
for systems are identified below with the compositional names \(EM\)
and \(EMC\). Their system-name roles are declared separately
from same-letter schema roles.

\subsection{The source language and the consequence convention}
\label{sec:p08-cn-language}

\begin{hbdefinition}{def:p08-cn-language}{Primitive possibility and the definitional bridge}
The 1980 source language has denumerably many atoms \(p_n\), truth
and falsity constants, the classical Boolean connectives, and primitive
\(\Box\) and \(\Diamond\). Its classical systems contain the
duality schema \(\Diamond A\leftrightarrow\neg\Box\neg A\).
The existing handbook sheet \ref{sheet:hilbert-e} uses \(\Box\)
and treats \(\Diamond\) as that defined abbreviation.
Let \(t(A)\) translate source formulas to this language: it fixes
atoms and constants, commutes with Boolean connectives and \(\Box\),
and sends \(\Diamond A\) to \(\neg\Box\neg t(A)\).
Source \(T\) and \(\bot\) translate to the handbook's truth and
falsity constants. This is a specified definitional extension, not an
identification of two undeclared primitive languages.

A source classical system \(L\) contains all propositional
tautologies and is closed under tautological consequence and
theorem-only replacement of equivalents. Write \(\vdash_L A\)
for theoremhood. Local \(\Gamma\vdash_L A\) uses finite
sequences whose lines are premises, theorems of \(L\), or
modus-ponens consequences of earlier lines. Applying modal replacement
requires a separate theorem derivation without \(\Gamma\), as in
\ref{rule:p08-hilbert-e-re}. The semantic local consequence requires
all premises at one world, not at every world in the model.
The source's Definition~2.14 states local deducibility through a
theorem implication from finitely many premises; the finite-sequence
formulation is proved equivalent in the deduction argument below.
\end{hbdefinition}

\begin{hbtheorem}{thm:p08-cn-language}{proved=annex:p8.cn.language}{Definitional equivalence and occurrence-specific replacement}
Every source classical system proves \(A\leftrightarrow t(A)\),
where the translated formula is read back as a source formula.
The primitive and defined presentations of E and its M/C/N extensions
therefore have the same translated theorems and local consequences.
Theorem-only replacement holds at any chosen occurrences inside a formula,
including under either modality. In particular theorem equivalence
\(A\leftrightarrow B\) gives both boxed and diamond equivalence;
it is not a rule for lifting merely local premise equivalence.
Equivalently, the source classical condition can use
\(\Box A\leftrightarrow\neg\Diamond\neg A\) and
theorem-only diamond replacement. Either duality schema together
with occurrence-specific replacement characterizes the same
classical closure.
\end{hbtheorem}

\begin{hbdefinition}{def:p08-cn-dual}{The source's formula and modality duals}
The source dual \(A^*\) negates atoms, swaps truth and falsity,
commutes with negation, swaps conjunction and disjunction, and swaps
\(\Box\) and \(\Diamond\). Its implication clause is
\((A\to B)^*=\neg A^*\land B^*\); its equivalence clause is
\((A\leftrightarrow B)^*=A^*\leftrightarrow\neg B^*\).
For a word of modal/negation operators, \(\mu^*\) swaps the two
modalities and leaves negation unchanged. Iteration \(\Box^0\)
or \(\Diamond^0\) means the identity; positive powers repeat
the stated operator.
Two modality words are equivalent in a system when their
corresponding formulas are theorem-equivalent at every formula argument.
\end{hbdefinition}

\begin{hblemma}{thm:p08-cn-dual}{proved=annex:p8.cn.dual}{Duality without a normality assumption}
Every classical system proves \(A^*\leftrightarrow\neg A\)
and \(A^{**}\leftrightarrow A\). It also proves
\(\mu A\leftrightarrow\neg\mu^*\neg A\).
Consequently \(\mu A\to\nu A\) is a theorem for every A
exactly when \(\nu^*A\to\mu^*A\) is a theorem for every A.
Equivalent modality words have equivalent dual words.
These are Chellas's replacement/duality principles for classical systems,
and require no K axiom or necessitation from premises
\cite[theorems~8.2--8.6, pp.~232--234]{Chellas1980}.
\end{hblemma}

\subsection{Minimal models, truth sets and closure operations}
\label{sec:p08-cn-models}

\begin{hbdefinition}{def:p08-cn-models}{The source model and its powerset operator}
The source minimal model is \(\mathcal M=(W,N,P)\), with
\(P_n\subseteq W\) and \(N(w)\subseteq\mathcal P(W)\).
It is the model of \ref{def:p08-neighborhood-model}, with
\(P_n=V(p_n)\). Write \(\|A\|_{\mathcal M}\) for its truth
set. Then
\[
 w\models\Box A\ \Longleftrightarrow\ \|A\|\in N(w),
 \qquad
 w\models\Diamond A\ \Longleftrightarrow\ W\setminus\|A\|\notin N(w).
\]
The source permits an empty \(W\). Such a model imposes no
constraint on global validity and supplies no pointed satisfiability
witness; excluding it therefore changes none of the theoremhood or
pointed consequences here. Every countermodel constructed below is nonempty.

Define the corresponding powerset operation
\(f(X)=\{w:X\in N(w)\}\). Necessity acts by \(f\) and
possibility by \(d(X)=W\setminus f(W\setminus X)\).
Conversely any operation \(f:\mathcal P(W)\to\mathcal P(W)\)
determines \(N(w)=\{X:w\in f(X)\}\).
\end{hbdefinition}

\begin{hblemma}{thm:p08-cn-semantics}{proved=annex:p8.cn.semantics}{Truth-set induction and extensional soundness}
The model and powerset-operation presentations agree at every world
for every formula. Their Boolean truth sets are complement, intersection
and union, with implication \((W\setminus\|A\|)\cup\|B\|\).
The duality schema is valid, and validity of \(A\leftrightarrow B\)
throughout a model implies validity of
\(\Box A\leftrightarrow\Box B\) throughout that model.
Truth depends only on the formula's occurring atoms when \(W,N\)
are fixed. These are extensional statements; equality of truth values
at one world is insufficient for modal replacement.
\end{hblemma}

\begin{hbdefinition}{def:p08-cn-closures}{Supplementation, intersections, quasi-filters and filters}
The source's supplemented condition is equivalently either
\[
 X\cap Y\in N(w)\ \Longrightarrow\ X,Y\in N(w)
 \quad\hbox{or}\quad
 X\in N(w),\ X\subseteq Y\ \Longrightarrow\ Y\in N(w).
\]
Its intersection condition uses binary intersections, or equivalently
all nonempty finite intersections. The unit condition is \(W\in N(w)\).
A quasi-filter has supplementation and intersection closure; a filter
additionally has the unit. Empty neighborhood families are allowed in
quasi-filters; filters need not exclude the empty proposition.

For a family \(\mathcal B\subseteq\mathcal P(W)\), let
\(\mathcal B^+\) be its superset closure, and
\(\mathcal B^-\) its closure under nonempty finite intersections.
The positive integer requirement is essential: intersection of zero
sets would insert the unit even when N is absent.
Quasi-filtering is \(\mathcal B^\#=(\mathcal B^+)^-\).
Apply these operations separately to each world's family, retaining
the worlds and atom valuation.
\end{hbdefinition}

\begin{hbtheorem}{thm:p08-cn-closure}{proved=annex:p8.cn.closure}{Exact closure properties and M/C/N correspondence}
The two supplemented conditions are equivalent, and
\((\mathcal B^+)^-=(\mathcal B^-)^+\).
Supplementation of an intersection-closed family remains
intersection-closed; intersection closure of a supplemented family
remains supplemented. For a supplemented family, nonemptiness is
equivalent to containing \(W\).

For neighborhood frames with arbitrary valuations, validity of M,
C and N corresponds respectively to supplementation, binary
intersection closure and containing the unit. For one fixed valuation,
formula validity alone need not constrain nondefinable proposition sets.
These closure transformations do not generally preserve the original
model's truth of every formula
\cite[secs.~7.1--7.2, pp.~207--219]{Chellas1980}.
\end{hbtheorem}

\subsection{Axioms, rules and the eight distinct systems}
\label{sec:p08-cn-systems}

\begin{hbdefinition}{def:p08-cn-systems}{The compositional systems and source names}
For \(S\subseteq\{M,C,N\}\), let \(L_S\) be E plus the
indicated axiom schemas on the exact existing sheets:
\ref{sheet:hilbert-e}, \ref{sheet:hilbert-m},
\ref{sheet:hilbert-c}, \ref{sheet:hilbert-n} and \ref{sheet:hilbert-r}.
M is \(\Box(A\land B)\to(\Box A\land\Box B)\);
C is \((\Box A\land\Box B)\to\Box(A\land B)\);
N is \(\Box\topv\). In this paragraph M/C/N are schema names.

Chellas also names his smallest monotonic system \(M\), generated
by the theorem rule RM taking an implication to its boxed implication;
his smallest regular system \(R\) uses theorem rule RR taking
\((A\land B)\to C\) to
\((\Box A\land\Box B)\to\Box C\).
His normal rule RK takes a theorem finite-conjunction implication
to the corresponding boxed implication, including the zero-premise
case. These are rules on theorems, not on unrestricted local premises.
The source's R schema is binary meet equivalence; its K schema is
\(\Box(A\to B)\to(\Box A\to\Box B)\).
Possessing the K schema alone does not make a system the normal
logic K.
\end{hbdefinition}

\begin{hbtheorem}{thm:p08-cn-rules}{proved=annex:p8.cn.rules}{Exact equivalence of source rule and axiom presentations}
For classical systems, RM is equivalent to adding schema M.
Regularity is equivalent to adding M and C, and hence
Chellas's systems satisfy \(M=EM\), \(R=EMC\).
In a monotonic system, schema C and the K schema are equivalent.
Regular systems have the K schema without necessarily having N.
Adding N to a regular system gives the normal rule RK for every
finite conjunction, including zero antecedents, and
\(K=EMCN\). Conversely the usual K-schema/theorem-only
necessitation presentation, with the declared Boolean base and
duality convention, supplies M, C, N and RE.
In a classical system N is equivalent to theorem-only necessitation;
this does not supply the other normality conditions.
The corresponding diamond principles follow by duality, and the
rules and finite-meet/finite-join principles iterate to every finite
modal depth \cite[secs.~8.1--8.2, pp.~231--243]{Chellas1980}.
\end{hbtheorem}

\begin{hbtheorem}{thm:p08-cn-separation}{proved=annex:p8.cn.separation}{Finite witnesses and every required strict distinction}
The following three source models isolate the independent axioms:
\begin{enumerate}
\item One world with empty neighborhood family satisfies M and C
      and falsifies N.
\item Two worlds \(a,b\), each with family
      \(\{\{a\},\{b\},W\}\), satisfy M and N and falsify C
      at the valuations \(V(p)=\{a\},V(q)=\{b\}\).
\item Two worlds, each with family \(\{\varnothing,W\}\),
      satisfy C and N and falsify M at those same valuations.
\end{enumerate}
Consequently \(L_S\subseteq L_T\) exactly when \(S\subseteq T\):
all eight systems E, EM, EC, EN, EMC, EMN, ECN and EMCN are distinct.
This is their inclusion diagram, not the full lattice of modal logics.
The third frame falsifies the K schema with \(A=\bot\) and
\(B=p\), so EN is not K. The first shows EMC is not EMCN.
E, EM, EC and EMC have no theorems of the form \(\Box A\).

The second family also excludes \(\varnothing\), so
\(\Diamond\topv\) holds everywhere, yet both \(\Box p\)
and \(\Box\neg p\) hold. The source's consistency-of-obligation
principle cannot be inferred from permission of truth without the
needed intersection assumption
\cite[theorems~7.4, 9.2; exercises~7.2, 9.7]{Chellas1980}.
\end{hbtheorem}

\subsection{Proof sets and the different canonical models}
\label{sec:p08-cn-canonical}

\begin{hbdefinition}{def:p08-cn-proof-sets}{Local consistency, maximal worlds and proof sets}
For a consistent classical system \(L\), a set \(\Gamma\)
is consistent if \(\Gamma\nvdash_L\bot\).
Let \(W_L\) consist of its maximal consistent sets.
The source's proof set notation is
\[
                   |A|_L=\{u\in W_L:A\in u\}.
\]
These bars denote the published proposition set, not cardinality.
Let \(D_L=\{|A|_L:A\text{ a formula}\}\).
The smallest canonical neighborhoods are
\(N^{\min}(u)=\{|A|_L:\Box A\in u\}\).
A canonical neighborhood family may add any non-proof sets,
but cannot change membership of a proof set. The largest is
\(N^{\max}(u)=N^{\min}(u)\cup(\mathcal P(W_L)\setminus D_L)\).
All canonical valuations assign \(p_n\) its proof set.
\end{hbdefinition}

\begin{hblemma}{thm:p08-cn-maximal}{proved=annex:p8.cn.maximal}{Deduction, countable extension and unambiguous proof sets}
The local deduction theorem holds. Every consistent set in this
denumerable language extends to a maximal consistent set by deciding
an enumeration of formulas while preserving consistency.
Such sets contain all theorems, are closed under modus ponens, and
have the classical Boolean membership conditions, including exactly
one of \(A,\neg A\).
Moreover
\[
 |A|_L\subseteq|B|_L\ \Longleftrightarrow\ \vdash_L A\to B,
 \qquad
 |A|_L=|B|_L\ \Longleftrightarrow\ \vdash_L A\leftrightarrow B.
\]
Thus theorem-only RE makes canonical membership unambiguous even
when different formulas define the same proposition.
\end{hblemma}

\begin{hbtheorem}{thm:p08-cn-canonical}{proved=annex:p8.cn.canonical}{The source canonical truth lemma and the eight constructions}
Every canonical model just defined satisfies
\(u\models A\) exactly when \(A\in u\).
For \(L_S\), use the smallest model if M is absent and the
supplementation of its smallest model if M is present. This model
satisfies every required C or N condition as well. In particular
the five required canonical constructions are:
\[
 E:N^{\min};\quad EC:N^{\min};\quad EN:N^{\min};\quad
 EM:(N^{\min})^+;\quad EMC:(N^{\min})^+.
\]
Their respective closures follow from the exact axiom/rule hypotheses;
supplementation is proved to preserve every canonical proof-set
membership before its truth lemma is used
\cite[theorems~9.3--9.14, pp.~252--259]{Chellas1980}.
\end{hbtheorem}

\begin{hbtheorem}{thm:p08-cn-completeness}{proved=annex:p8.cn.completeness}{Theorem determination and separately proved strong local consequence}
For each of the eight \(L_S\), its theorems are exactly the
formulas valid on neighborhood frames with the corresponding M/C/N
conditions. More strongly, under the stated theorem-only premise
convention, for any set \(\Gamma\),
\[
 \Gamma\vdash_{L_S}A
 \quad\Longleftrightarrow\quad
 \text{every such pointed model satisfying }\Gamma\text{ satisfies }A.
\]
This strong local extension is proved here; the source's determination
statements are about theoremhood. Local semantic consequence is compact.
Lifting an equivalence obtained only from local premises is unsound
even for K: at a point seeing another world, two atoms can agree at
the first point and differ at its successor, giving unequal box truth.
\end{hbtheorem}

\subsection{Augmentation and the precise normal boundary}
\label{sec:p08-cn-augmentation}

\begin{hbdefinition}{def:p08-cn-augmentation}{Principal neighborhoods and augmentation}
A neighborhood frame is augmented if, at every world,
\[
 X\in N(w)\quad\Longleftrightarrow\quad\bigcap N(w)\subseteq X.
\]
This includes the requirement that its core \(\bigcap N(w)\)
belongs to \(N(w)\). Its augmentation replaces each family by
all supersets of that core, with \(\bigcap\varnothing=W\).
For an accessibility relation \(R\), its principal neighborhoods
are \(N_R(w)=\{X:R(w)\subseteq X\}\).
\end{hbdefinition}

\begin{hbtheorem}{thm:p08-cn-augmentation}{proved=annex:p8.cn.augmentation}{Relational equivalence, finite filters and canonical augmentation}
Every standard relational model is pointwise equivalent to its
principal-neighborhood model, and every augmented model to the
relational model with successors \(\bigcap N(w)\).
Every finite filter is augmented. An infinite filter need not be:
on \(\mathbb R\), the source family at \(a\) consists of all
sets containing some open interval \((a,b)\), \(b>a\).
It is a filter whose core is empty and whose family excludes the
empty set. It therefore has no same-world relational representation
preserving truth for every valuation.

For any consistent normal system, augmentation of the smallest
canonical model is still canonical. Its proof uses the theorem-only
finite boxed-consequence argument and an actual maximal-world witness
for each allegedly new proof-set membership. Hence K is determined
both by filters and by augmented frames; equality of the resulting
logic does not identify all infinite filters with relational frames
\cite[secs.~7.3, 9.3, pp.~220--223, 259--261]{Chellas1980}.
\end{hbtheorem}

\begin{hbtheorem}{thm:p08-cn-interval}{proved=annex:p8.cn.interval}{Exact intervals and local persistence are different models}
If \(N(a)\) consists only of open intervals containing \(a\),
then \(\Box A\) asserts that \(\|A\|\) is itself such
an interval. Its supplementation instead asserts that the truth set
contains an open interval around \(a\). At zero, the set
\((-1,1)\cup\{3\}\) distinguishes the two clauses.
The latter is the local-persistence reading; it must not be inferred
from membership in the former family without the supplementation.
\end{hbtheorem}

\subsection{Finite truth preservation and effective decisions}
\label{sec:p08-cn-finite}

\begin{hbdefinition}{def:p08-cn-filtration}{The source's finest finite observations}
Let \(F\) be a finite set of source formulas closed under
subformulas. If a unit-containing frame is required, include
\(\Box\topv\) and its subformula. Identify worlds agreeing on
every member of \(F\), and write \(\pi:W\to W_F\) for the
quotient map. Observed propositions are truth sets of \(A\)
when \(\Box A\in F\), and complements of truth sets of \(A\)
when \(\Diamond A\in F\). The finest filtered family at
\([w]\) contains exactly their images \(\pi[X]\) for which
\(X\in N(w)\). Atom valuations in \(F\) use their images;
unmentioned atoms can be assigned arbitrarily.
\end{hbdefinition}

\begin{hbtheorem}{thm:p08-cn-filtration}{proved=annex:p8.cn.filtration}{Finest filtrations, each closure case and the unit}
The filtered families are well defined on equivalence classes.
Every \(A\in F\) has the same truth at \(w\) and \([w]\).
If the original model is supplemented, supplement its finest
filtration; if it is intersection-closed, take nonempty finite
intersection closure; if it has both properties, quasi-filter it.
Each remains an \(F\)-filtration and has the required closure.
The added observation \(\Box\topv\) preserves the unit in
all four constructions. No generated neighborhood can create a
wrong box or diamond value on \(F\); the proof checks this
using preimages of quotient-observed sets in each case.
There are at most \(2^{|F|}\) quotient worlds
\cite[theorems~7.12--7.17; exercise~7.49]{Chellas1980}.
\end{hbtheorem}

\begin{hbtheorem}{thm:p08-cn-decision}{proved=annex:p8.cn.decision}{Finite-model property and an actual terminating decision procedure}
All eight M/C/N combinations, hence E, EM, EC, EN and EMC, are
determined by their finite neighborhood models. For a formula \(A\),
take its subformulas and, when N is required, those of
\(\Box\topv\). A failed theorem has a countermodel of at most
\(2^{|F|}\) worlds. Enumerating all neighborhood tables and
valuations on up to that many worlds, retaining the required closures,
and evaluating \(A\) terminates and decides theoremhood.
It returns a concrete pointed countermodel if theoremhood fails.

The source's second route also works: finite schemas with effective
theorem-only RPL/RE give a positive proof search, and finite model
enumeration gives a negative search; dovetail them. Completeness and
finite determination ensure one terminates. Finite-premise local
consequence is decided by its deduction-theorem implication.
No sharper complexity bound or validity of an infinite-premise
decision algorithm is asserted
\cite[sec.~2.8; theorems~9.22--9.28; exercises~9.40, 9.44]{Chellas1980}.
\end{hbtheorem}

\subsection{Algebraic forms and a finite modal projection}
\label{sec:p08-cn-algebra}

\begin{hbdefinition}{def:p08-cn-algebra}{Modal Boolean algebras and the logical quotient}
A modal Boolean algebra is a Boolean algebra \(B\) with an
arbitrary unary operation \(\ast:B\to B\).
Interpret necessity by \(\ast\), possibility by \(\neg\ast\neg\),
and Boolean formulas by their Boolean operations. A formula is true
when its value is the unit. The M/C/N conditions are respectively
\[
 \ast(a\land b)\leq\ast a\land\ast b,\qquad
 \ast a\land\ast b\leq\ast(a\land b),\qquad \ast1=1.
\]
The algebraic elements are propositions, not necessarily worlds.
In the logical quotient identify formulas when their equivalence
is a theorem of \(L\); write their classes \([A]_L\).
The Boolean operations are induced by syntax, and the modal operation
is \(\ast[A]_L=[\Box A]_L\). This equivalence class differs
from the source proof set \(|A|_L\).
\end{hbdefinition}

\begin{hbtheorem}{thm:p08-cn-algebra}{proved=annex:p8.cn.algebra}{Well-defined algebraic semantics and determination}
The powerset operation model is an algebraic model and has the
same formula truth: all-world truth means value \(W\).
The logical quotient is well defined, including its modal operation
by theorem-only RE, and its canonical evaluation sends \(A\)
to \([A]_L\). This value is the unit exactly when \(\vdash_L A\).
Each of the eight systems is therefore determined by modal Boolean
algebras with its respective M/C/N conditions.
\end{hbtheorem}

\begin{hbtheorem}{thm:p08-cn-algebra-finite}{proved=annex:p8.cn.algebra-finite}{The finite Boolean projection that preserves relevant modal values}
In an algebraic countermodel for \(A\), let \(B_F\) be the
finite Boolean subalgebra generated by the values of its subformulas.
For \(b\in B\), define
\(r(b)=\bigvee\{c\in B_F:c\leq b\}\), and on \(B_F\) put
\(\ast_Fc=r(\ast c)\). This retains the value of every subformula,
and retains each required M/C/N condition. It gives a finite
algebraic countermodel for each failed theorem of the eight systems.
Agreement with the original modal operation is needed on relevant
arguments whose modal outputs are represented; it need not hold on
every element of the Boolean subalgebra. Modal output outside that
subalgebra cannot be its new operation's value.
\end{hbtheorem}

\subsection{Other source conditions and illustrative canonical choices}
\label{sec:p08-cn-other}

\begin{hbdefinition}{def:p08-cn-other-conditions}{The source's D/T/B/4/5 conditions}
For \(X\subseteq W\), write \(f(X)=\{v:X\in N(v)\}\)
and \(d(X)=W\setminus f(W\setminus X)\) as above.
The corresponding frame conditions are:
\begin{enumerate}
\item[(d)] \(X\in N(w)\) implies \(W\setminus X\notin N(w)\).
\item[(t)] \(X\in N(w)\) implies \(w\in X\).
\item[(b)] \(w\in X\) implies \(d(X)\in N(w)\).
\item[(iv)] \(X\in N(w)\) implies \(f(X)\in N(w)\).
\item[(v)] \(X\notin N(w)\) implies \(W\setminus f(X)\in N(w)\).
\end{enumerate}
These are conditions on proposition families. They are not
automatically seriality, reflexivity, symmetry, transitivity or
Euclideanness of an undeclared accessibility relation.
\end{hbdefinition}

\begin{hbtheorem}{thm:p08-cn-other-conditions}{proved=annex:p8.cn.other-conditions}{Schema validation and iterated source conditions}
With arbitrary valuations, these conditions correspond respectively
to the source schemas D, T, B, 4 and 5. Their diamond duals follow
by the proved duality. Iterated box and diamond truth sets are
\(f^n(X)\) and \(W\setminus f^n(W\setminus X)\).
For nonnegative integers \(k,l,m,n\), the general source schema
\(\Diamond^k\Box^l A\to\Box^m\Diamond^n A\)
corresponds to
\[
 W\setminus f^k(W\setminus f^l(X))
       \subseteq f^m(W\setminus f^n(W\setminus X)).
\]
Its converse reverses that inclusion. The particular \(G\)
case has all four indices one. This establishes the needed
operator conditions without assuming normality
\cite[sec.~7.4, pp.~223--226]{Chellas1980}.
\end{hbtheorem}

\begin{hbtheorem}{thm:p08-cn-t5-principles}{proved=annex:p8.cn.t5-principles}{The source ET5 principles and exact six modalities}
Every classical system containing T and 5 has N,
\(\Diamond\topv\), theorem-only box and diamond necessitation,
and D, B and 4. It has the reductions
\[
 \begin{aligned}
 \Box\Box A&\leftrightarrow\Box A,&
 \Box\Diamond A&\leftrightarrow\Diamond A,\\
 \Diamond\Box A&\leftrightarrow\Box A,&
 \Diamond\Diamond A&\leftrightarrow\Diamond A.
 \end{aligned}
\]
The smallest such system ET5 has exactly the six modalities
identity, box, diamond and their negations; among the affirmative
three, box implies identity and identity implies diamond, with
neither reverse implication valid. ET5 does not contain M and is
not normal: a finite neighborhood counterexample is given in the annex
\cite[theorems~8.15--8.16; exercise~9.9]{Chellas1980}.
\end{hbtheorem}

\begin{hbtheorem}{thm:p08-cn-t5-completeness}{proved=annex:p8.cn.t5-completeness}{ET, E5 and ET5 use different canonical choices}
ET is determined by frames satisfying (t), using its smallest
canonical model. E5 is determined by frames satisfying (v), using
its largest canonical model. ET5 is determined by frames satisfying
both, using its smallest canonical model and the derived N theorem.
Their source theoremhood determination extends to strong local
consequence under the fixed premise convention by the supplied
maximal-set argument. These results do not assert that arbitrary
combinations of additional axioms have the same canonical choice
\cite[theorems~9.17--9.21, pp.~262--263]{Chellas1980}.
\end{hbtheorem}

\begin{hbtheorem}{thm:p08-cn-nonproof}{proved=annex:p8.cn.nonproof}{The restriction belongs to omitted non-proof sets}
For a canonical model of a classical E5-extension, condition (v)
is automatic on proof sets. On other sets the remaining requirement
is: whenever a non-proof set \(X\) is omitted from \(N(w)\),
its negative-membership extension \(\{v:X\notin N(v)\}\)
must belong to \(N(w)\). The largest canonical choice satisfies
this because it omits no non-proof sets. Deleting them all is not
generally safe for E5 alone: there is a maximal world without
\(\Box\topv\), while every non-proof set is absent at all
worlds of its smallest model. This produces a failure of (v).
The existence of such a set and world is constructed in the annex.
\end{hbtheorem}

The source setup and semantic constructions above are attributed to
this 1980 presentation. Scott's generalized framework, Montague's
\emph{Pragmatics and Intensional Logic}, and Pacuit's later development
remain separately identified sources. The proof of a stated naming
or semantic bridge does not erase their different original presentations.
